\section{Calcul discret de la mesure sur la limite projective : différences, cocycle de retenue, ultramétrique, reconstruction de \texorpdfstring{\(\mathbb R\)}{R} et désintégration généalogique}
\label{sec:complejo-discreto-medida-rev2}
\index{mesure!complexe discret}
\index{chemin!intégration}
\index{retenue!cocycle}

La distinction entre compter et mesurer admet une formulation opératoire. Compter
assigne un cardinal à une collection de marques. Mesurer évalue des relations
entre marques au moyen d'une règle locale, les intègre le long d'un chemin et
conserve la carte dans laquelle le résultat est publié. L'unité minimale de ce
processus est une différence ; sa composition produit une longueur, une circulation et,
lorsqu'une réduction modulaire intervient, une retenue. Cette architecture situe la
mesure avant sa coordonnée réelle : le scalaire publié est la contraction d'un
calcul qui possède un domaine, une orientation et une mémoire.

\subsection{Marques et opérateur de différences}

Soit \(P_n\) le chemin aux sommets ordonnés
\(V(P_n)=\{v_0,\ldots,v_{n-1}\}\) et aux arêtes orientées
\(e_j=(v_{j-1},v_j)\), \(1\leq j\leq n-1\). Pour un anneau commutatif
\(A\), nous définissons
\[
 C^0(P_n;A)=A^{V(P_n)},
 \qquad
 C^1(P_n;A)=A^{E(P_n)},
\]
et l'opérateur de différences
\begin{equation}
 d:C^0(P_n;A)\longrightarrow C^1(P_n;A),
 \qquad
 (df)(e_j)=f(v_j)-f(v_{j-1}).
 \label{eq:diferencia-cero-uno-rev2}
\end{equation}

\begin{theorem}[Rang de l'opérateur de différences sur un chemin]
\label{thm:rango-diferencias-camino-rev2}
Si \(A\) est un corps, alors
\[
 \ker d=A\mathbf 1,
 \qquad
 \operatorname{rank}d=n-1.
\]
Par conséquent, \(n\) marques reliées déterminent \(n-1\) différences
indépendantes et un unique mode constant.
\end{theorem}

\begin{proof}
L'égalité \(df=0\) implique successivement
\(f(v_0)=f(v_1)=\cdots=f(v_{n-1})\), de sorte que le noyau est formé des
fonctions constantes et a pour dimension un. Le théorème du
rang donne \(\operatorname{rank}d=n-1\). De manière équivalente, toute
assignation \((a_1,\ldots,a_{n-1})\) aux arêtes s'intègre de manière unique
une fois \(f(v_0)\) fixé :
\[
 f(v_k)=f(v_0)+\sum_{j=1}^{k}a_j.
\]
\end{proof}

Le mode constant exprime la liberté de choisir une origine. Les différences
contiennent l'information relationnelle ; la marque initiale fournit la section
qui reconstruit les positions. Dans la cellule inaugurale de HMT, neuf marques
ordonnées produisent huit différences intérieures. La clôture périodique ajoute la
neuvième arête et transforme le chemin en cycle ; la complétion bidimensionnelle
de ces intervalles sera réalisée de manière exacte dans la construction d'APP.

\subsection{Longueur positive et intégrale orientée}

Deux évaluations d'une route doivent rester séparées. Une métrique discrète
est une cochaîne positive
\[
 \ell:E(P_n)\longrightarrow\mathbb R_{>0}.
\]
Pour un chemin non orienté \(|\gamma|\), sa longueur est
\begin{equation}
 L_\ell(\gamma)=\sum_{e\in|\gamma|}m_\gamma(e)\,\ell(e),
 \label{eq:longitud-positiva-camino-rev2}
\end{equation}
où \(m_\gamma(e)\) compte les traversées de l'arête. Une cochaîne
orientée \(\omega\in C^1(P_n;A)\), en revanche, s'intègre avec un signe :
\begin{equation}
 \int_\gamma\omega
 :=\sum_{k=1}^{r}\varepsilon_k\omega(e_{i_k}),
 \qquad \varepsilon_k\in\{-1,+1\}.
 \label{eq:integral-cocadena-camino-rev2}
\end{equation}
La longueur somme le coût géométrique ; l'intégrale orientée somme le transport. La
première est invariante par inversion de la route et la seconde change de signe.

\begin{proposition}[Théorème fondamental discret]
\label{prop:teorema-fundamental-discreto-rev2}
Pour \(f\in C^0(P_n;A)\) et toute route orientée \(\gamma:a\to b\),
\[
 \int_\gamma df=f(b)-f(a).
\]
En particulier, l'intégrale d'une cochaîne exacte sur un cycle est nulle.
\end{proposition}

\begin{proof}
La somme des différences consécutives est télescopique. Chaque sommet intérieur
apparaît une fois avec un signe positif et une fois avec un signe négatif ; seules
les extrémités subsistent.
\end{proof}

La proposition distingue deux niveaux de clôture. Le groupe
\(H_1(C_n;\mathbb Z)\cong\mathbb Z\) fournit le support topologique du
tour ; une holonomie requiert en outre un transport. Si
\(T:\operatorname{Path}(C_n)\to\mathcal C\) est un foncteur vers une catégorie
d'automorphismes, l'holonomie d'un cycle \(\gamma\) est
\begin{equation}
 \operatorname{Hol}_T(\gamma)=T(e_r)\cdots T(e_1).
 \label{eq:holonomia-transporte-ciclo-rev2}
\end{equation}
Le cycle apporte la classe de parcours ; \(T\) apporte la connexion discrète qui
transforme cette classe en un endomorphisme. Cette distinction sera essentielle lorsque
le retour visible retrouvera la phase et que le transport conservera un état de
frontière différent.

\subsection{Retenue comme cocycle de composition}

La réduction modulaire publie une classe et une section choisit des représentants.
Soit \(Q=\mathbb Z/9\mathbb Z\) et soit
\(s_0:Q\to\{0,\ldots,8\}\subset\mathbb Z\) la section résiduelle ordinaire.
Nous définissons
\begin{equation}
 \kappa_9(a,b)
 :=\frac{s_0(a)+s_0(b)-s_0(a+b)}{9}\in\mathbb Z.
 \label{eq:cociclo-acarreo-nonadico-rev2}
\end{equation}
Alors
\[
 s_0(a)+s_0(b)=s_0(a+b)+9\kappa_9(a,b).
\]

\begin{proposition}[Identité cocyclique de la retenue]
\label{prop:acarreo-dos-cociclo-rev2}
L’application \(\kappa_9:Q\times Q\to\mathbb Z\) vérifie
\begin{equation}
 \kappa_9(a,b)+\kappa_9(a+b,c)
 =\kappa_9(b,c)+\kappa_9(a,b+c).
 \label{eq:identidad-dos-cociclo-rev2}
\end{equation}
C'est donc un \(2\)-cocycle normalisé associé à la section choisie.
\end{proposition}

\begin{proof}
On développe les deux membres au moyen de
\eqref{eq:cociclo-acarreo-nonadico-rev2}. Après multiplication par neuf, les deux
sont égaux à
\[
 s_0(a)+s_0(b)+s_0(c)-s_0(a+b+c).
\]
La normalisation s'obtient à partir de \(s_0(0)=0\).
\end{proof}

L'identité cocyclique exprime l'associativité du relèvement : regrouper d'abord
\((a+b)+c\) ou \(a+(b+c)\) produit la même classe visible et la même retenue
totale. Dans une catégorie de routes, la même structure prend la forme
\[
 \widetilde g\,\widetilde h
 =\kappa(g,h)\,\widetilde{gh},
\]
où \(\widetilde g,\widetilde h\) sont des transports enrichis qui relèvent des
morphismes visibles. La retenue est ainsi un défaut de composition parfaitement
contrôlé avant de recevoir une quelconque interprétation géométrique comme
courbure ou torsion.

\subsection{Compatibilité des cartes nonadique, ternaire et de complétion décimale en base \texorpdfstring{\(10^3\)}{10^3}}
\label{subsec:compatibilidad-lectores-9-3-1000-rev3}

Les lecteurs résiduels qui interviennent dans TPK ne sont pas des noms différents d'
une même opération. Soient
\[
 R_9:=\mathbb Z/9\mathbb Z,
 \qquad
 r_9:\mathbb Z\longrightarrow R_9,
 \qquad
 r_3:\mathbb Z\longrightarrow\mathbb F_3
\]
les réductions canoniques, et soit
\[
 \pi_3:R_9\longrightarrow\mathbb F_3,
 \qquad \pi_3(\overline n)=n\bmod3.
\]
Le carré des réductions commute :
\begin{equation}
\begin{CD}
 \mathbb Z @>{r_9}>> R_9\\
 @V{r_3}VV @VV{\pi_3}V\\
 \mathbb F_3 @= \mathbb F_3 .
\end{CD}
\qquad
 r_3=\pi_3\circ r_9.
\label{eq:diagrama-reducciones-9-3-rev3}
\end{equation}
La section positive
\[
 s_+:R_9\longrightarrow\{1,\ldots,9\},
 \qquad
 s_+(\overline0)=9,
 \qquad
 s_+(\overline k)=k\quad(1\le k\le8),
\]
produit, pour \(n>0\), la racine numérique positive
\[
 \operatorname{dr}_9^+(n)=s_+(r_9(n)).
\]
Cette section n'est pas un homomorphisme : elle publie la classe nulle par \(9\) et
oublie le quotient. La section résiduelle \(s_0(\overline0)=0\) publie cette
même classe par \(0\). Par conséquent, le changement visible \(9\mapsto0\) est un
changement de section, et l'égalité
\[
 9=s_0(\overline0)+9\cdot1
\]
conserve comme quotient la retenue \(1\). Le successeur topologique du cycle
positif est, en revanche, \(9\mapsto1\), avec
\[
 10=s_+(\overline1)+9\cdot1.
\]
Ainsi se séparent la clôture \(P_9\to C_9\), le changement de carte et l'incrément
de mémoire ; aucun des trois ne se substitue aux deux autres.

La coordonnée du radical appartient à un autre diagramme. Pour
\(\mathfrak m=(3)\subset R_9\), l'application
\[
 \iota_{\mathfrak m}:\mathbb F_3\longrightarrow\mathfrak m,
 \qquad \overline a\longmapsto3a,
\]
est un isomorphisme d'espaces vectoriels sur \(\mathbb F_3\), d'inverse
\begin{equation}
 \eta_{\mathfrak m}:\mathfrak m\longrightarrow\mathbb F_3,
 \qquad
 \eta_{\mathfrak m}(3a)=a\bmod3.
 \label{eq:coordenada-radical-compatibilidad-rev3}
\end{equation}
Ce n'est pas un isomorphisme d'anneaux. Elle ne coïncide pas non plus avec la réduction directe :
pour \(3a\in\mathfrak m\), on a \(\pi_3(3a)=0\), tandis que
\(\eta_{\mathfrak m}(3a)=a\bmod3\). La première lecture publie la classe
résiduelle ; la seconde publie la coordonnée interne du radical.

Enfin, soient \(u_j\in\{0,\ldots,999\}\) et
\[
 D_N=1000D_{N-1}+u_N,
 \qquad D_0=0,
\]
la concaténation de blocs en base \(10^3\). Comme \(10^3\equiv1\pmod9\),
\begin{equation}
 r_9(D_N)=r_9\!\left(\sum_{j=1}^{N}u_j\right),
 \qquad
 r_3(D_N)=r_3\!\left(\sum_{j=1}^{N}u_j\right).
 \label{eq:compatibilidad-base-mil-reducciones-rev3}
\end{equation}
La base \(10^3\) conserve l'ordre des blocs ; les réductions ne conservent que
leur classe accumulée. La racine numérique ajoute une section positive ; la coordonnée
radicale retrouve une direction interne de \((3)\). Cette hiérarchie fixe quelle
information chaque lecteur publie et quel quotient, ordre ou carte la
mémoire TPK doit retenir pour reconstruire l'antécédent.

\subsection{Système gradué de mesure et registre non contracté}

Les trois niveaux précédents s'organisent en un système gradué :
\[
 \boxed{
 \begin{array}{ccl}
 \text{degré }0&:&\text{valeurs ou états sur des marques},\\
 \text{degré }1&:&\text{différences et transports sur les arêtes},\\
 \text{degré }2&:&\text{cocycles de composition et de retenue}.
 \end{array}}
\]
Les degrés zéro et un forment le complexe cellulaire
\(C^0\xrightarrow d C^1\) ; le degré deux vit dans le complexe du nerf de la
catégorie des chemins. Leur réunion constitue le \emph{complexe discret de
mesure} : une structure organisatrice qui conserve la différence entre
géométrie cellulaire et mémoire catégorielle.

Pour une route enrichie
\(\Gamma=(x_0\xrightarrow{e_1}x_1\to\cdots\xrightarrow{e_r}x_r)\), le
registre est l'intégrale non contractée
\begin{equation}
 \Lambda(\Gamma)
 =\bigl((x_{k-1},e_k,x_k),\omega_k,\kappa_k,\sigma_k,b_k,m_k\bigr)_{k=1}^{r},
 \label{eq:ledger-integral-no-contraida-rev2}
\end{equation}
où \(\omega_k\) est la lecture locale, \(\kappa_k\) la retenue, \(\sigma_k\)
l'orientation ou signature partielle, \(b_k\) la donnée de frontière et \(m_k\) la
mémoire conservée. Un observable est une application typée
\[
 \mathcal O=F_{\mathcal O}\circ\Lambda.
\]
L'évaluation peut publier une somme, une classe résiduelle, une signature ou un
scalaire sans transformer ces sorties en états complets. La concaténation des
routes concatène les registres et ajoute le cocycle à la frontière commune ; le
registre non contracté conserve donc exactement l'information nécessaire pour recomposer l'
histoire qu'une coordonnée isolée contracte.

\begin{figure}[tbp]
\centering
\begin{tikzpicture}[
  >=Latex,
  every node/.style={font=\small,align=center},
  object/.style={draw,rounded corners=2pt,inner sep=5pt,text width=3.00cm,
                 minimum height=1.25cm},
  arrow/.style={->,semithick},
  maplabel/.style={font=\footnotesize,fill=white,inner sep=1.5pt}
]
\node[object] (marks) at (0,0)
  {marques et états\\\(C^0\)};
\node[object] (edges) at (5.25,0)
  {différences et transport\\\(C^1\)};
\node[object] (paths) at (10.50,0)
  {chemins et cycles\\intégration et holonomie};
\node[object] (memory) at (0,-2.35)
  {cocycle de composition\\retenue et frontière};
\node[object] (section) at (5.25,-2.35)
  {registre non contracté\\section compatible};
\node[object] (outputs) at (10.50,-2.35)
  {évaluations typées\\valeur, signature ou symétrie};
\draw[arrow] (marks) -- node[maplabel,above]{\(d\)} (edges);
\draw[arrow] (edges) -- node[maplabel,above]{\(\int_\gamma\)} (paths);
\draw[arrow] (paths) -- node[maplabel,right]{transporte TPK} (outputs);
\draw[arrow] (memory) -- node[maplabel,above]{composition} (section);
\draw[arrow] (section) -- node[maplabel,above]{contraction} (outputs);
\draw[arrow] (edges) -- node[maplabel,left]{défaut de relèvement} (memory);
\draw[arrow] (paths) -- node[maplabel,right]{histoire} (section);
\end{tikzpicture}
\caption{Architecture du complexe discret de mesure. L'opérateur de
différences, l'intégrale de chemins et le cocycle de composition occupent des
domaines distincts ; le TPK les réunit dans une section compatible avant qu'
un lecteur ne publie une valeur ou une symétrie.}
\label{fig:complejo-discreto-medida-rev3}
\end{figure}

\subsection{Reconstruction du continu à partir du système projectif de cylindres compatibles}

Le prolongement à une profondeur arbitraire se formule au moyen d'un système
projectif. Soient \(S_n\) des ensembles finis non vides de préfixes enrichis et
\(r_{n+1,n}:S_{n+1}\to S_n\) des applications de restriction surjectives. Leur
espace de sections compatibles est
\[
 \Omega_\infty=\varprojlim_n S_n
 =\{(x_n)_n:r_{n+1,n}(x_{n+1})=x_n\}.
\]
Chaque \(S_n\) est muni de la topologie discrète et \(\Omega_\infty\) de la
topologie de sous-espace du produit.

\subsection{Cylindres, ultramétrique et retour archimédien avec conservation de la mémoire}

Pour \(a\in S_n\), le cylindre de préfixe \(a\) est
\begin{equation}
 [a]_n:=\{x\in\Omega_\infty:x_n=a\}.
 \label{eq:cilindro-genealogico-rev3}
\end{equation}
Deux cylindres d'un même niveau sont disjoints, et chaque cylindre de niveau
\(n+1\) est contenu dans le cylindre de sa restriction. Cette base d'
ouverts et fermés publie la profondeur partagée par deux sections sans
contracter leur mémoire en une coordonnée réelle.

Pour \(x\neq y\), soit
\[
 N(x,y):=\min\{n\geq0:x_n\neq y_n\}.
\]
Fijado \(0<\vartheta<1\), definimos
\begin{equation}
 d_\vartheta(x,y):=\vartheta^{N(x,y)},
 \qquad d_\vartheta(x,x):=0.
 \label{eq:ultrametrica-genealogica-rev3}
\end{equation}

\begin{proposition}[Ultramétrique de première divergence]
\label{prop:ultrametrica-genealogica-rev3}
La fonction \(d_\vartheta\) est une ultramétrique et ses boules sont précisément
les cylindres, avec le réajustement d'indice correspondant au niveau initial :
\[
 d_\vartheta(x,z)
 \leq
 \max\{d_\vartheta(x,y),d_\vartheta(y,z)\}.
\]
Par conséquent, deux sections qui publient la même valeur peuvent rester à
distance positive lorsque divergent leur feuillet, leur retenue, leur frontière ou leur
histoire de prolongement.
\end{proposition}

\begin{proof}
Si \(x\) et \(y\) coïncident jusqu'au niveau \(n\), ainsi que \(y\) et \(z\),
alors \(x\) et \(z\) coïncident jusqu'à ce niveau. La première divergence du
premier et du troisième état ne se produit donc pas avant le minimum des
deux autres. Comme \(0<\vartheta<1\), élever \(\vartheta\) à ces
exposants produit l'inégalité forte. La description des boules découle
directement de la fixation d'un préfixe commun.
\end{proof}

\begin{theorem}[Compacité et complétude de l'espace généalogique]
\label{thm:compacidad-ultrametrica-genealogica-rev3}
Si chaque \(S_n\) est fini et si les applications de liaison sont surjectives, alors
\(\Omega_\infty\), muni de \(d_\vartheta\), est non vide, compact et complet.
Si de plus
\begin{equation}
 \forall x\in\Omega_\infty\ \forall n\ \exists y\neq x:
 y_n=x_n,
 \label{eq:perfeccion-genealogica-rev3}
\end{equation}
il ne possède aucun point isolé. Dans ce régime, c'est un compact métrisable,
de dimension zéro et parfait ; par la caractérisation classique, il est homéomorphe
à l'espace de Cantor.
\end{theorem}

\begin{proof}
La surjectivité et la ramification finie fournissent une section par le
lemme de König. Le produit des ensembles finis discrets est compact, et
les équations \(r_{n+1,n}(x_{n+1})=x_n\) définissent un sous-espace fermé. La
topologie induite coïncide avec la topologie des cylindres et, par la
\cref{prop:ultrametrica-genealogica-rev3}, avec celle de
\(d_\vartheta\). Tout espace métrique compact est complet. La condition
\eqref{eq:perfeccion-genealogica-rev3} garantit que chaque cylindre contient
une section distincte de son centre, de sorte qu'aucun singleton n'est ouvert.
\end{proof}

La périodicité visible et le retour de l'état appartiennent à des applications
distinctes. Soient \(p:E\twoheadrightarrow B\) une projection, \(T:E\to E\) un
transport enrichi et \(f:B\to B\) l'évolution publiée, avec
\(p\circ T=f\circ p\). Si \(f^9=I_B\), alors
\begin{equation}
 p\circ T^9=p,
 \label{eq:retorno-visible-fibra-rev3}
\end{equation}
de sorte que \(T^9\) préserve chaque fibre de \(p\), sans nécessairement être
l'identité sur celle-ci. Dans la coordonnée nonadique
\[
 (q,r)\longmapsto
 \begin{cases}
  (q,r+1),&r<9,\\
  (q+1,1),&r=9,
 \end{cases}
\]
un tour retrouve la position et augmente la hauteur. Dans la périodicité
CT108, neuf itérations font avancer d'une phase dodécaphasique et cent huit
retrouvent la position et la phase ; l'état enrichi peut encore conserver
des feuillets, des frontières, des orientations et des généalogies différents. C'est la distinction
formelle entre retour de phase et retour d'état.

\begin{theorem}[Reconstruction de la valeur par sections compatibles]
\label{thm:reconstruccion-continuo-secciones-rev2}
L'espace \(\Omega_\infty\) est non vide et compact. Supposons qu'à chaque
préfixe \(x_n\in S_n\) soit assigné un intervalle compact
\(I_n(x_n)\subset\mathbb R\) de sorte que, pour toute section
\(x=(x_n)_n\),
\[
 I_{n+1}(x_{n+1})\subseteq I_n(x_n),
 \qquad
 \operatorname{diam}I_n(x_n)\longrightarrow0.
\]
Alors il existe une unique valeur
\begin{equation}
 \operatorname{ev}_{\mathbb R}(x)
 \in\bigcap_{n\geq0}I_n(x_n).
 \label{eq:evaluacion-seccion-continua-rev2}
\end{equation}
Si les assignations d'intervalles sont localement constantes à chaque niveau, l'
évaluation \(\operatorname{ev}_{\mathbb R}:\Omega_\infty\to\mathbb R\) est
continue.
\end{theorem}

\begin{proof}
La surjectivité permet de construire par récurrence des préfixes compatibles ; la
compacité des ensembles finis et la propriété d'intersection finie donnent
une section globale. Le produit \(\prod_nS_n\) est compact et
\(\Omega_\infty\) est fermé, car chaque équation de compatibilité définit un
fermé. Pour une section fixée, les intervalles sont compacts, emboîtés et de
diamètre tendant vers zéro ; le théorème des intervalles emboîtés fournit un
unique point commun. Étant donné \(\varepsilon>0\), on choisit un niveau dont l'intervalle
a un diamètre inférieur à \(\varepsilon\). Les sections qui partagent ce
préfixe ont leurs évaluations dans le même intervalle, ce qui prouve la
continuité.
\end{proof}

La droite réelle apparaît ainsi comme codomaine d'une contraction. Deux sections
peuvent partager une valeur et conserver des registres différents ; la relation
\[
 x\sim_{\mathbb R}y
 \quad\Longleftrightarrow\quad
 \operatorname{ev}_{\mathbb R}(x)=\operatorname{ev}_{\mathbb R}(y)
\]
forme le quotient archimédien. La généalogie demeure dans
\(\Omega_\infty\), tandis que la coordonnée vit dans son image réelle. Le problème
du continu reçoit donc une réponse constructive dans le domaine
couvert par les systèmes de préfixes déclarés : la continuité de l'
évaluation et la mémoire de la section sont des propriétés coordonnées, non des
descriptions rivales.

\subsection{Couverture cohérente et reconstruction du quotient archimédien}

L'évaluation précédente assigne une valeur à chaque section. Pour démontrer que l'
espace visible est entièrement reconstruit, il faut en outre que le
raffinement ne perde aucun point. Soit \(K\subset\mathbb R\) compact et supposons
que les intervalles du théorème précédent satisfassent
\begin{align}
 K&=\bigcup_{a\in S_0} I_0(a),
 \label{eq:cobertura-inicial-continuo-rev3}\\
 I_n(a)&=\bigcup_{\substack{b\in S_{n+1}\\r_{n+1,n}(b)=a}}I_{n+1}(b)
 \qquad(a\in S_n).
 \label{eq:cobertura-hijos-continuo-rev3}
\end{align}
Ces identités définissent une \emph{couverture cohérente} : chaque compact
visible se décompose exactement en les compacts de ses prolongements.

\begin{theorem}[Reconstruction généalogique d'un compact]
\label{thm:cociente-genealogico-compacto-rev3}
Supposons que :
\begin{enumerate}
\item les ensembles \(S_n\) sont finis et non vides, et les applications de liaison
      \(r_{n+1,n}\) sont surjectives ;
\item les intervalles sont compacts non vides, emboîtés et uniformément
      contractants ;
\item les conditions
      \eqref{eq:cobertura-inicial-continuo-rev3} et
      \eqref{eq:cobertura-hijos-continuo-rev3} sont satisfaites.
\end{enumerate}
Alors \(\operatorname{ev}_{\mathbb R}(\Omega_\infty)=K\), et l'application
induite
\begin{equation}
 \overline{\operatorname{ev}}_{\mathbb R}:
 \Omega_\infty/\!\sim_{\mathbb R}\xrightarrow{\ \cong\ }K
 \label{eq:homeomorfismo-cociente-continuo-rev3}
\end{equation}
est un homéomorphisme. La fibre au-dessus de \(y\in K\) est l'ensemble complet des
sections compatibles qui publient cette même valeur.
\end{theorem}

\begin{proof}
La compacité et la continuité proviennent du
\cref{thm:reconstruccion-continuo-secciones-rev2}. \(y\in K\) étant fixé, soit
\[
 T_n(y):=\{a\in S_n:y\in I_n(a)\}.
\]
La couverture cohérente rend tous les \(T_n(y)\) non vides et relie chaque
état à un prolongement. L'arbre obtenu est à ramification finie et
possède des sommets à toute profondeur ; par le lemme de la branche infinie, il admet une
chaîne compatible \((x_n)_n\). Le point \(y\) appartient à tous les
\(I_n(x_n)\), et la contraction impose
\(\operatorname{ev}_{\mathbb R}(x)=y\). Par conséquent, l'évaluation est
surjective. L'application du quotient est continue et bijective ; comme son
domaine est compact et que \(K\) est séparé, c'est un homéomorphisme.
\end{proof}

Le théorème distingue deux affirmations. La contraction produit une valeur unique
pour chaque généalogie ; la couverture cohérente démontre que toutes les valeurs
du compact apparaissent. Une vérification à profondeur finie certifie les
niveaux examinés, mais ne remplace aucune de ces hypothèses infinies.

\begin{proposition}[Atlas généalogique de la droite réelle]
\label{prop:atlas-genealogico-recta-real-rev3}
Pour chaque \(m\geq1\), soit \(K_m=[-m,m]\) et supposons donné un système qui
satisfait le \cref{thm:cociente-genealogico-compacto-rev3}, avec des
inclusions fermées
\[
 \jmath_m:\Omega_\infty^{(m)}\hookrightarrow\Omega_\infty^{(m+1)}
\]
qui entrelacent les évaluations et l'inclusion
\(K_m\hookrightarrow K_{m+1}\). Alors le quotient archimédien de la
colimite stricte des sections est
\[
 \varinjlim_m K_m=\bigcup_{m\geq1}[-m,m]=\mathbb R,
\]
muni de sa topologie usuelle, tandis que les fibres de la colimite conservent les
généalogies de chaque valeur.
\end{proposition}

\begin{proof}
Chaque morceau possède le quotient \(K_m\) par
\eqref{eq:homeomorfismo-cociente-continuo-rev3}. Les inclusions déclarées
rendent ces quotients compatibles. La topologie finale de l'exhaustion par
intervalles compacts coïncide avec la topologie usuelle de \(\mathbb R\) : un
ensemble est ouvert exactement lorsque son intersection avec chaque \(K_m\) est
ouverte dans \(K_m\).
\end{proof}

\subsection{Descente algébrique des opérations}

Reconstruire l'espace visible ne suffit pas à reconstruire son algèbre. Soit
\(\star\) une opération partielle continue sur \(\Omega_\infty\), de domaine
\(D_\star\) saturé par les fibres de
\(\operatorname{ev}_{\mathbb R}\times\operatorname{ev}_{\mathbb R}\).

\begin{theorem}[Critère de descente par fibres]
\label{thm:descenso-operaciones-continuo-rev3}
Il existe une opération visible \(\overline\star\) qui satisfait
\[
 \operatorname{ev}_{\mathbb R}(x\star y)
 =\operatorname{ev}_{\mathbb R}(x)\,\overline\star\,
  \operatorname{ev}_{\mathbb R}(y)
\]
si et seulement si la valeur du membre de gauche est constante lorsqu'on remplace
\(x\) ou \(y\) par tout autre représentant de leurs fibres respectives.
Si deux opérations continues descendues coïncident avec la somme et le produit
sur des sous-ensembles denses de leurs domaines visibles, elles coïncident avec eux sur tout
le domaine.
\end{theorem}

\begin{proof}
La constance sur les fibres est exactement la propriété universelle du
quotient : elle permet de définir \(\overline\star\) indépendamment du représentant,
et elle est nécessaire pour que cette définition ait un sens. La dernière affirmation
découle de la continuité et de la densité.
\end{proof}

APP conserve la donnée requise pour aborder cette descente : les feuillets somme et
produit agissent sur le même support, mais gardent des quotients, des retenues et des
fibres différents. L'opération visible peut descendre bien que la mémoire de
ses représentants demeure distincte.

\subsection{Mesure projective et conservation de l'unité}

La topologie des cylindres organise la proximité et la troncature ; une théorie des
chemins requiert en outre des poids. Soit \(\mu_n\) une probabilité sur \(S_n\).
La famille est projectivement cohérente lorsque
\begin{equation}
 \mu_n(a)=
 \sum_{\substack{b\in S_{n+1}\\r_{n+1,n}(b)=a}}\mu_{n+1}(b).
 \label{eq:consistencia-medida-cilindrica-rev3}
\end{equation}
L'égalité exprime la conservation locale de l'unité : le poids d'une
histoire tronquée est la somme des poids de tous ses prolongements.

\begin{theorem}[Mesure généalogique sur la limite projective]
\label{thm:medida-genealogica-limite-rev3}
Toute famille de probabilités finies qui satisfait
\eqref{eq:consistencia-medida-cilindrica-rev3} détermine une unique
probabilité borélienne \(\mu_\infty\) sur \(\Omega_\infty\) telle que
\begin{equation}
 \mu_\infty([a]_n)=\mu_n(a)
 \label{eq:medida-cilindro-rev3}
\end{equation}
pour tout cylindre de préfixe \(a\in S_n\).
\end{theorem}

\begin{proof}
Les unions finies de cylindres forment une algèbre. Chacune peut
être raffinée en une union disjointe à un même niveau, et
\eqref{eq:consistencia-medida-cilindrica-rev3} démontre que la somme de ses
poids ne dépend pas du niveau choisi. On obtient ainsi une prémesure. Comme les
cylindres sont ouverts et fermés et que \(\Omega_\infty\) est compact, une
suite décroissante d'éléments de l'algèbre d'intersection vide atteint
déjà le vide à un certain niveau. La prémesure est donc continue au vide
et s'étend de manière unique à la \(\sigma\)-algèbre borélienne.
\end{proof}

L'évaluation transporte cette mesure vers le compact visible,
\begin{equation}
 \nu=(\operatorname{ev}_{\mathbb R})_*\mu_\infty,
 \qquad
 \nu(B)=\mu_\infty\bigl(\operatorname{ev}_{\mathbb R}^{-1}(B)\bigr).
 \label{eq:pushforward-medida-visible-rev3}
\end{equation}
La probabilité visible agrège les généalogies d'une même fibre sans
les effacer de l'espace total.

\begin{proposition}[Désintégration par valeurs et multiplicité généalogique]
\label{prop:desintegracion-valores-genealogicos-rev3}
Si \(K\) est un espace borélien standard, il existe, pour \(\nu\)-presque tout
\(y\in K\), une probabilité conditionnelle \(\mu_y\) portée par
\[
 \mathfrak G_y:=\operatorname{ev}_{\mathbb R}^{-1}(y)
\]
telle que, pour toute fonction intégrable \(F\),
\begin{equation}
 \int_{\Omega_\infty}F\,\mathrm d\mu_\infty
 =\int_K
   \left(\int_{\mathfrak G_y}F\,\mathrm d\mu_y\right)
   \mathrm d\nu(y).
 \label{eq:desintegracion-fibras-visibles-rev3}
\end{equation}
Ainsi, la coordonnée visible et la distribution des histoires qui la réalisent
sont des données typées distinctes : un observable qui se factorise par
\(\operatorname{ev}_{\mathbb R}\) ne lit que \(\nu\), tandis qu'un lecteur de
mémoire peut distinguer les mesures \(\mu_y\) au sein d'une même fibre.
\end{proposition}

\begin{proposition}[Conditionnement compatible des cylindres]
\label{prop:condicionamiento-cilindrico-compatible-rev3}
Soit \(A\subseteq\Omega_\infty\) une union de cylindres de niveau \(m\), avec
\(\mu_\infty(A)>0\). Pour tout \(n\geq m\), la distribution conditionnelle
sur \(S_n\) s'obtient en supprimant les états dont la troncature n'
appartient pas à \(A\) et en renormalisant les poids restants. Les distributions
obtenues satisfont toujours la cohérence projective
\eqref{eq:consistencia-medida-cilindrica-rev3}.
\end{proposition}

\begin{proof}
C'est la règle de Bayes sur les ensembles finis de niveau \(n\). La somme des
poids conditionnels des prolongements coïncide avec le poids conditionnel du
parent parce que l'intersection avec \(A\) se distribue sur l'union disjointe de
ses prolongements. Le conditionnement élimine les futurs incompatibles et
redistribue l'unité entre les survivants ; il ne crée pas de nouvelles histoires et n'
efface pas de l'espace total celles restées en dehors de l'événement lu.
\end{proof}

Cette séparation fournit la base exacte permettant de distinguer la préparation,
la lecture, le conditionnement et l'histoire dans la réalisation quantique.

\begin{statusbox}[title={Force mathématique et fonction architecturale}]
Le rang \(n-1\), le théorème fondamental discret, l'identité cocyclique de la
retenue, les théorèmes de reconstruction et la mesure projective sont des résultats
exacts dans les domaines déclarés. Leur application HMT utilise comme structure de départ les états,
les restrictions, les intervalles et les transports construits par APP--TRIT--TPK. Le
complexe discret de mesure est une formalisation organisatrice : il n'identifie pas
le cocycle catégoriel de degré deux à une cellule géométrique tant que l'
application correspondante n'est pas fournie. Cette séparation permet au langage
commun de gouverner ensuite l'holonomie, la somme sur les chemins et la transition du
nombre à la particule sans anticiper leurs réalisations physiques.
\end{statusbox}
